\section{Discussion}

\subsection{Stale reads and timeouts}

The matched histories show two different outcomes under the partition. With
weaker settings, an operation could complete using stale state from the
isolated branch, producing a consistency violation. With stronger settings,
the operation could instead time out when the required state was unavailable.

The three configuration controls contributed differently to these outcomes.
Read concern determined whether a read could return the isolated member's
local state or was restricted to majority-committed data. Write concern
determined whether a write could be acknowledged on the isolated branch
without reaching a majority. When causal consistency was enabled, the session
also carried the logical dependency from an earlier operation to the next one.
As illustrated by the C6 RYW history in
Figure~\ref{fig:representative-trace}, the dependent read therefore waited for
the required state and eventually timed out rather than returning an older
value.

\subsection{Consistency, completion, and latency}

Table~\ref{tab:completion-partition-outcomes} reports the decidable-history
rate $D$ and operation completion rate $C$ separately.
Figure~\ref{fig:decidability-latency} compares decidability and
critical-operation latency for the selected C1 and C6 partition experiments.

\begin{table}[ht]
\centering
\small
\setlength{\tabcolsep}{3pt}
\caption{History outcomes, decidability, critical-operation responses, and latency under the recorded partition.}
\label{tab:completion-partition-outcomes}
\begin{tabular}{@{}l l r l r l r@{}}
\toprule
Config. & Property & $n_{\text{valid}}$ & P/V/U/I & $D$ & \shortstack{$C$\\completed/issued (rate)} & All-attempt p95 (ms) \\
\midrule
\RQFourPartitionRows
\bottomrule
\end{tabular}
\end{table}

\maybefigure[fig:decidability-latency]{submission/figures/rq4_decidability_latency.pdf}{Paired C1/C6 RYW and MW cells during the recorded partition. The upper panel shows the decidable-history rate $D$; the lower panel gives all-attempt p95 latency of the critical operation. Operation completion rate $C$ is reported in Table~\ref{tab:completion-partition-outcomes}. ``Not issued'' marks a cell without a critical-operation record. Each cell contains 20 valid histories.}


The contrast between C1 and C6 shows why consistency results should be
interpreted together with completion. Under C1, all 20 RYW and MW histories
were decidable and exposed violations, with the critical operations completing
in every trial. Under C6, none of the corresponding histories was decidable,
and the critical operations were not issued because an earlier operation had
already failed to complete or left the required evidence unresolved. Thus, the
absence of an observed C6 violation does not mean that these trials demonstrated
the consistency property; instead, the partition prevented a complete history
from being obtained.

Latency provides another view of this trade-off. Operations that exposed stale
or reordered state under C1 completed quickly. Under stronger settings, an
earlier operation could instead wait until its timeout when the required state
or acknowledgement could not be obtained. These earlier-operation timeouts are
therefore absent from the critical-operation latency in
Figure~\ref{fig:decidability-latency}, because the critical operations were
never issued. Their delay is instead reflected in the fault-level all-attempt
latency reported in Table~\ref{tab:fault-summary}.
Appendix~\ref{app:completion-details} separately reports all-attempt and
resolved critical-operation latency for the selected configuration histories.

\FloatBarrier